\RequirePackage{plautopatch}
\documentclass[platex,dvipdfmx,a4paper]{jsarticle}

\usepackage{amsmath}
\usepackage{amsfonts}
\usepackage{amssymb}

\usepackage[dvipdfmx]{graphicx}
\usepackage{float}

\usepackage{tikz}
\usepackage{caption}

% 「%」は以降の内容を「改行コードも含めて」無視するコマンド
\usepackage[%
  dvipdfmx,% 欧文ではコメントアウトする
  setpagesize=false,%
  bookmarks=true,%
  bookmarksdepth=tocdepth,%
  bookmarksnumbered=true,%
  colorlinks=false,%
  pdftitle={},%
  pdfsubject={},%
  pdfauthor={},%
  pdfkeywords={}%
]{hyperref}
% PDFのしおり機能の日本語文字化けを防ぐ((u)pLaTeXのときのみかく)
\usepackage{pxjahyper}

\usepackage{framed}
\setlength{\OuterFrameSep}{0pt}

\newcommand{\hb}[2]{\frac{\partial #1}{\partial #2}}
\newcommand{\bbn}[2]{\frac{d #1}{d #2}}
\newcommand{\nhb}[3]{\frac{\partial^{#1} #2}{\partial #3^{#1}}}
\newcommand{\nbbn}[3]{\frac{d^{#1} #2}{d #3^{#1}}}
\newcommand{\vectxy}[2]{\begin{bmatrix}#1\\#2\end{bmatrix}}
\newcommand{\vectxyz}[3]{\begin{bmatrix}#1\\#2\\#3\end{bmatrix}}
\newcommand{\vect}[1]{\boldsymbol{\mathrm{#1}}}

\newcommand{\ssum}[3]{\sum_{#1=#2}^#3}

\newcommand{\veff}{V_{\mathrm{eff}}(r)}

\hypersetup{colorlinks=true}
\hypersetup{linkcolor=red}

\title{解析力学演習レポート2}
\author{25B00729　嶋田 悠音}


\begin{document}
\maketitle
\clearpage

\section*{アンテナ線・オーディオ接続コードに表示される$\Omega$の意味}

\subsection*{1．はじめに}

同軸ケーブルや平行フィーダー線には、$50\,\Omega$、$75\,\Omega$、$300\,\Omega$などの表示がある。また、オーディオ機器にも、入力インピーダンス、出力インピーダンス、スピーカーの公称インピーダンスとして$\Omega$が表示されている。

これらはいずれも電圧と電流の比という意味ではインピーダンスであるが、同じ$\Omega$表記でも表している物理量は必ずしも同じではない。特に、ケーブルの特性インピーダンス、機器の入出力インピーダンス、ケーブル導体の抵抗を区別する必要がある。

アンテナ線に表示される$50\,\Omega$、$75\,\Omega$、$300\,\Omega$などは、通常はケーブル全体の直流抵抗ではなく、伝送線路としての特性インピーダンスを表している。一方、オーディオ機器の端子やスピーカーに表示される$\Omega$は、機器または負荷のインピーダンスを表す場合が多い。そこで、まず電信方程式から特性インピーダンスを導き、その後、アンテナ線とオーディオ接続における意味の違いを考える。

\subsection*{2．分布定数回路としてのケーブル}

通常の低周波回路では、抵抗器、コンデンサー、コイルなどの回路素子が特定の場所に集中していると考える。このような回路を集中定数回路という。

しかし、十分に長いケーブルや高周波信号を扱うケーブルでは、導体の抵抗、自己インダクタンス、導体間の静電容量、絶縁体を通じた漏れがケーブル全体に連続的に分布している。そのため、ケーブルを次の単位長さあたりの量で表す：
\begin{itemize}
  \item $R$：単位長さあたりの直列抵抗 $\mathrm{[\Omega/m]}$
  \item $L$：単位長さあたりの直列インダクタンス $\mathrm{[H/m]}$
  \item $G$：単位長さあたりの並列コンダクタンス $\mathrm{[S/m]}$
  \item $C$：単位長さあたりの並列静電容量 $\mathrm{[F/m]}$
\end{itemize}
このように、回路定数が導線に沿って連続的に分布している回路を分布定数回路という。

ケーブル上の位置を$x$とし、その位置における電圧と電流をそれぞれ$V(x,t)$、$I(x,t)$とする。長さ$\Delta x$の微小区間にKirchhoffの電圧則を適用する：
\begin{align*}
  V(x,t)-V(x+\Delta x,t)
  &=R\Delta x I(x,t)+L\Delta x\hb{I(x,t)}{t}。
\end{align*}
両辺を$\Delta x$で割り、$\Delta x\to0$の極限を取る：
\begin{align*}
  \hb{V}{x}=-RI-L\hb{I}{t}。
\end{align*}
同様に、微小区間にKirchhoffの電流則を適用し、$\Delta x\to0$の極限を取る：
\begin{align*}
  \hb{I}{x}=-GV-C\hb{V}{t}。
\end{align*}
したがって、分布定数回路における電圧と電流は、次の電信方程式を満たす：
\begin{align*}
  \hb{V}{x}
  &=-RI-L\hb{I}{t}, \\
  \hb{I}{x}
  &=-GV-C\hb{V}{t}。
\end{align*}
この方程式は、ケーブル中の電圧と電流が場所によらず同時に変化するのではなく、波として有限の速度で伝わることを表している。

\subsection*{3．電信方程式と特性インピーダンス}

角周波数$\omega$の正弦波信号を考え、電圧と電流をフェーザ表示する。時間依存性を$e^{j\omega t}$とすれば、電信方程式は次のようになる：
\begin{align*}
  \bbn{V}{x}
  &=-(R+j\omega L)I, \\
  \bbn{I}{x}
  &=-(G+j\omega C)V。
\end{align*}
第1式をさらに$x$で微分し、第2式を代入する：
\begin{align*}
  \nbbn{2}{V}{x}
  &=-(R+j\omega L)\bbn{I}{x} \\
  &=(R+j\omega L)(G+j\omega C)V。
\end{align*}
同様に、電流についても次式を得る：
\begin{align*}
  \nbbn{2}{I}{x}
  &=(R+j\omega L)(G+j\omega C)I。
\end{align*}
ここで、伝搬定数$\gamma$を次のように定義する：
\begin{align*}
  \gamma
  &=\sqrt{(R+j\omega L)(G+j\omega C)}
   =\alpha+j\beta。
\end{align*}
ただし、$\alpha$は減衰定数、$\beta$は位相定数である。このとき、電圧と電流の一般解は、正の$x$方向へ進む波と負の$x$方向へ進む波の重ね合わせとして、次のように表される：
\begin{align*}
  V(x)
  &=V^+e^{-\gamma x}+V^-e^{\gamma x}。
\end{align*}
これを$x$で微分して電信方程式へ代入する：
\begin{align*}
  I(x)
  &=\frac{\gamma}{R+j\omega L}
    \left(V^+e^{-\gamma x}-V^-e^{\gamma x}\right)。
\end{align*}
ここで、特性インピーダンス$Z_0$を次のように定義する：
\begin{align*}
  Z_0
  &=\frac{R+j\omega L}{\gamma}
   =\sqrt{\frac{R+j\omega L}{G+j\omega C}}。
\end{align*}
この定義を用いれば、電流は次のように書ける：
\begin{align*}
  I(x)
  &=\frac{V^+}{Z_0}e^{-\gamma x}
    -\frac{V^-}{Z_0}e^{\gamma x}。
\end{align*}
したがって、正の$x$方向へ進む波については、電圧と電流の比が次のようになる：
\begin{align*}
  \frac{V^+e^{-\gamma x}}{(V^+/Z_0)e^{-\gamma x}}
  =Z_0。
\end{align*}
すなわち、特性インピーダンスは、一方向へ進む電圧波と電流波の比である。

ケーブルの損失を無視できる場合には$R=G=0$とおけるので、特性インピーダンスは次のように簡単になる：
\begin{align*}
  Z_0
  &=\sqrt{\frac{j\omega L}{j\omega C}}
   =\sqrt{\frac{L}{C}}。
\end{align*}
また、このとき伝搬定数は次式で与えられる：
\begin{align*}
  \gamma
  &=\sqrt{(j\omega L)(j\omega C)}
   =j\omega\sqrt{LC}。
\end{align*}
よって、位相定数は$\beta=\omega\sqrt{LC}$であり、伝搬速度$v$は次のようになる：
\begin{align*}
  v
  &=\frac{\omega}{\beta}
   =\frac{1}{\sqrt{LC}}。
\end{align*}

ここで重要なのは、損失のない線路でも$Z_0$は有限の値をもつことである。したがって、例えば「$75\,\Omega$ケーブル」とは、導体が$75\,\Omega$の直流抵抗をもち、その中で電力が熱として消費されるという意味ではない。特性インピーダンスは、電磁波として進む電圧波と電流波の比を表す量である。

また、特性インピーダンスは主としてケーブルの断面形状、導体間隔、導体径、絶縁体の誘電率などによって決まり、理想的に一様なケーブルであればその長さには依存しない。したがって、$1\,\mathrm{m}$の$75\,\Omega$ケーブルと$100\,\mathrm{m}$の$75\,\Omega$ケーブルでは、導体抵抗や減衰量は異なるが、特性インピーダンスはともに$75\,\Omega$である。

\subsection*{4．終端で生じる反射}

伝送線路の終端に負荷インピーダンス$Z_L$を接続する。簡単のため、負荷の位置を$x=0$とする。この位置における電圧と電流は、進行波と反射波を用いて次のように表される：
\begin{align*}
  V(0)
  &=V^++V^-, \\
  I(0)
  &=\frac{V^+}{Z_0}-\frac{V^-}{Z_0}。
\end{align*}
終端では$Z_L=V(0)/I(0)$が成り立つので、次式を得る：
\begin{align*}
  Z_L
  &=Z_0\frac{V^++V^-}{V^+-V^-}。
\end{align*}
ここで、電圧反射係数$\Gamma$を次のように定義する：
\begin{align*}
  \Gamma
  &=\frac{V^-}{V^+}。
\end{align*}
先の式を$\Gamma$について解く：
\begin{align*}
  Z_L(1-\Gamma)
  &=Z_0(1+\Gamma), \\
  \Gamma
  &=\frac{Z_L-Z_0}{Z_L+Z_0}。
\end{align*}
負荷インピーダンスが特性インピーダンスと等しい場合には、$Z_L=Z_0$であるから、$\Gamma=0$となる。したがって反射波は生じず、入射したエネルギーは負荷へ受け渡される。

一方、$Z_L\neq Z_0$の場合には反射が生じる。終端が開放されて$Z_L\to\infty$であれば$\Gamma=1$となり、終端が短絡されて$Z_L=0$であれば$\Gamma=-1$となる。したがって、どちらの場合も反射波の振幅は入射波の振幅と等しくなり、入射波は全反射される。ただし、短絡端では反射電圧の符号が反転する。

以上より、特性インピーダンスは、ケーブルの終端に接続したときに反射を生じさせない負荷インピーダンスであるとも解釈できる。また、無限に長い線路では終端に到達した波が戻ってこないため、その入力インピーダンスは$Z_0$に等しい。

\subsection*{5．アンテナ線における$\Omega$表示の意味}

アンテナ線では、用途に応じて異なる特性インピーダンスが用いられる。代表的なものとして、無線送信機などに用いられる$50\,\Omega$の同軸ケーブル、テレビ受信などに用いられる$75\,\Omega$の同軸ケーブル、$300\,\Omega$程度の平行フィーダー線などがある。これらの数値は、いずれもアンテナ線の特性インピーダンスを表す。

アンテナから高周波信号を受信機へ送る場合には、アンテナの給電点インピーダンス、ケーブルの特性インピーダンス、受信機の入力インピーダンスを同じ規格にそろえることが望ましい。例えば、$75\,\Omega$系では次の関係が成立するように構成する：
\begin{align*}
  Z_{\mathrm{antenna}}
  \simeq Z_0
  \simeq Z_{\mathrm{receiver}}
  \simeq75\,\Omega。
\end{align*}
送信の場合にも、送信機、給電線、アンテナの間でインピーダンスが大きく異なると反射が生じる。反射された波は進行波と重なり、給電線上に定在波を作る。この結果、アンテナへ供給される電力が減少し、反射が大きい場合には送信機の出力回路に過大な電圧または電流が生じる可能性もある。

ただし、アンテナ自体のインピーダンスは、一般には抵抗成分だけでなくリアクタンス成分も含む複素数である。したがって、「$50\,\Omega$アンテナ」という表現は、アンテナの直流抵抗が$50\,\Omega$であることを意味するのではなく、使用周波数付近で給電点インピーダンスが$50\,\Omega$系に整合するよう設計されていることを意味する。

\subsection*{6．伝送線路として扱う必要が生じる条件}

ケーブルを伝送線路として扱う必要があるかどうかは、ケーブル長と信号の波長との関係によって決まる。正弦波の場合、ケーブル内の波長$\lambda$は次式で与えられる：
\begin{align*}
  \lambda
  &=\frac{v}{f}。
\end{align*}
ここで、$v$はケーブル中の伝搬速度であり、一般には真空中の光速$c$より小さい。例えば、伝搬速度を$v=2.0\times10^8\,\mathrm{m/s}$とする。このとき、$f=20\,\mathrm{kHz}$の信号に対する波長は次のようになる：
\begin{align*}
  \lambda
  &=\frac{2.0\times10^8}{2.0\times10^4}
   =1.0\times10^4\,\mathrm{m}
   =10\,\mathrm{km}。
\end{align*}
数$\mathrm{m}$程度のケーブルは、この波長に比べて極めて短い。したがって、$20\,\mathrm{Hz}$から$20\,\mathrm{kHz}$程度のアナログ音声信号では、通常の長さのケーブルを伝送線路として扱う必要はほとんどない。この場合、ケーブル全体で電圧と電流がほぼ同時に変化するとみなし、集中定数回路として扱うことができる。

ただし、デジタル信号や矩形波については、単にパルスの繰返し周波数だけで判断することはできない。矩形波の立ち上がりと立ち下がりには高周波成分が含まれるため、信号の繰返し周波数が低くても、立ち上がり時間が短ければ伝送線路として扱う必要がある。

\subsection*{7．アナログオーディオ機器の接続}

\subsubsection*{7.1．ライン接続における入出力インピーダンス}

アナログオーディオ機器のライン出力とライン入力を接続する場合、アンテナ線の場合のように入出力インピーダンスを等しくすることは一般的ではない。通常は次の関係を満たすように設計する：
\begin{align*}
  Z_{\mathrm{out}}\ll Z_{\mathrm{in}}。
\end{align*}
このような接続をブリッジ接続または電圧伝送という。

出力機器を信号電圧$V_s$と出力インピーダンス$Z_{\mathrm{out}}$の直列回路で表し、入力機器を入力インピーダンス$Z_{\mathrm{in}}$で表す。電圧分割則を用いると、入力機器に加わる電圧は次式で与えられる：
\begin{align*}
  V_{\mathrm{in}}
  &=V_s\frac{Z_{\mathrm{in}}}{Z_{\mathrm{out}}+Z_{\mathrm{in}}}。
\end{align*}
$Z_{\mathrm{in}}\gg Z_{\mathrm{out}}$であれば、入力電圧は次のように近似できる：
\begin{align*}
  V_{\mathrm{in}}
  &\simeq V_s。
\end{align*}
したがって、信号電圧をほとんど低下させずに次の機器へ送ることができる。

例えば、出力インピーダンスが$100\,\Omega$、入力インピーダンスが$10\,\mathrm{k\Omega}$の場合を考える。このとき、入力電圧と信号電圧の比は次のようになる：
\begin{align*}
  \frac{V_{\mathrm{in}}}{V_s}
  &=\frac{10000}{100+10000}
   \simeq0.990。
\end{align*}
したがって、信号電圧の約$99\,\%$が入力機器に加わる。よって、アナログオーディオでは、「出力が$100\,\Omega$なら入力も$100\,\Omega$にする」という意味での整合は通常行わない。これは、高周波伝送線路において反射を防ぐためのインピーダンス整合とは目的が異なる。

\subsubsection*{7.2．ケーブルの静電容量の影響}

アナログオーディオケーブルでは、特性インピーダンスよりも静電容量が問題になる場合がある。ケーブル全体の静電容量を$C_{\mathrm{cable}}$、機器の出力抵抗を$R_{\mathrm{out}}$とすると、これらは近似的にRC低域通過回路を作る。この回路の遮断周波数$f_c$は次式で与えられる：
\begin{align*}
  f_c
  &=\frac{1}{2\pi R_{\mathrm{out}}C_{\mathrm{cable}}}。
\end{align*}
したがって、出力インピーダンスが高い機器に長くて静電容量の大きいケーブルを接続すると、可聴周波数帯の高周波成分が減衰する可能性がある。特に、受動型エレクトリックギターのピックアップのように出力インピーダンスが高い信号源では、ケーブル容量が周波数特性に大きく影響する。一方、通常のライン出力は比較的低い出力インピーダンスをもつため、数$\mathrm{m}$程度の接続ではケーブル容量の影響は小さい。

\subsection*{8．スピーカー接続における$\Omega$表示}

スピーカーに表示される$4\,\Omega$、$6\,\Omega$、$8\,\Omega$などは、ケーブルの特性インピーダンスではなく、スピーカーの公称負荷インピーダンスである。

スピーカーは、ボイスコイルの抵抗とインダクタンス、振動板の機械的な共振、エンクロージャーの音響的性質などを含む。そのため、そのインピーダンスは周波数によって変化する。したがって、「$8\,\Omega$」と表示されていても、すべての周波数で正確に$8\,\Omega$になるわけではない。

例えば、$8\,\Omega$のスピーカーへ実効値$20\,\mathrm{V}$を加え、スピーカーを純抵抗で近似する。このとき、電流と電力は次のようになる：
\begin{align*}
  I
  &=\frac{20}{8}
   =2.5\,\mathrm{A}, \\
  P
  &=\frac{20^2}{8}
   =50\,\mathrm{W}。
\end{align*}
このため、スピーカーケーブルでは特性インピーダンスを$8\,\Omega$にするのではなく、導体抵抗をスピーカーのインピーダンスに比べて十分小さくすることが重要である。ケーブル抵抗が大きいと、ケーブル自身で電圧降下と電力損失が生じ、スピーカーへ加わる電圧が低下する。

したがって、「$8\,\Omega$スピーカーには$8\,\Omega$の特性インピーダンスをもつケーブルが必要である」という理解は誤りである。必要なのは、アンプが$8\,\Omega$程度の負荷を安全に駆動できることと、十分に太く短いケーブルを用いて導体抵抗を小さくすることである。

\subsection*{9．デジタルオーディオ接続}

デジタルオーディオ接続では、音声の最高周波数が$20\,\mathrm{kHz}$程度であっても、ケーブルを流れる電気信号は高速なデジタルパルスである。そのため、伝送線路としての特性インピーダンスが重要になる。

代表的な規格として、民生用の同軸S/PDIFでは特性インピーダンスが$75\,\Omega$の同軸ケーブルが用いられる。また、業務用のAES3では、一般に特性インピーダンスが$110\,\Omega$の平衡ツイストペアケーブルが用いられる。

この場合には、送信回路、ケーブル、コネクター、受信回路を規定のインピーダンス系にそろえる必要がある。途中でインピーダンスが不連続になると、その位置で反射が発生する。反射波が入射波と重なることで、デジタル波形にはリンギングや立ち上がり部分の乱れが生じる。

小さな反射がただちに聴感上の音質変化になるとは限らない。しかし、波形の乱れが受信回路の許容範囲を超えれば、クロック抽出の不安定化、ビット判定の誤り、ロック外れ、クリックノイズ、音切れなどの原因になり得る。

外形が同じRCAケーブルでも、アナログ音声用ケーブルと$75\,\Omega$同軸デジタルケーブルは電気的仕様が同じとは限らない。短距離では通常のRCAケーブルでも動作する場合があるが、それはインピーダンスが適合していることを保証するものではない。したがって、安定した接続のためには、同軸S/PDIFには規格に従って$75\,\Omega$の同軸ケーブルを使用することが適切である。

\subsection*{10．インピーダンス整合の意味}

インピーダンス整合には、異なる二つの文脈がある。一つは、伝送線路において負荷インピーダンスを特性インピーダンスと一致させ、次の条件によって信号の反射を防ぐ整合である：
\begin{align*}
  Z_L=Z_0。
\end{align*}
アンテナ線、映像信号線、デジタル音声線などでは、こちらが重要である。

もう一つは、集中定数回路において負荷へ最大電力を供給するための共役整合である。信号源の内部インピーダンスが$Z_s$であるとき、負荷へ供給される平均電力が最大となる条件は次式で与えられる：
\begin{align*}
  Z_L=Z_s^*。
\end{align*}
ここで、$Z_s^*$は$Z_s$の複素共役である。

しかし、通常のアナログオーディオのライン接続では、最大電力を負荷へ送ることが目的ではなく、信号電圧を正確に取り出すことが目的である。そのため、次の条件を用いる：
\begin{align*}
  Z_{\mathrm{in}}\gg Z_{\mathrm{out}}。
\end{align*}
したがって、「インピーダンスは常に同じ値にそろえるべきである」という理解は正しくない。何を伝送するか、ケーブルが電気的に長いか、電圧と電力のどちらを重視するかによって、適切な接続方法は異なる。

\subsection*{11．結論}

同軸ケーブルや平行フィーダー線に表示される$50\,\Omega$、$75\,\Omega$、$300\,\Omega$などの値は、通常はケーブルの特性インピーダンスを表している。これは直流抵抗ではなく、分布定数回路を記述する電信方程式から得られる、一方向へ進む電圧波と電流波の比である。

負荷インピーダンスが特性インピーダンスと等しければ反射は生じない。したがって、アンテナ線では、送受信機、給電線、アンテナを同じインピーダンス系に合わせることが重要である。

一方、アナログオーディオのライン接続では、通常のケーブル長は音声信号の波長に比べて極めて短い。そのため、ケーブルの特性インピーダンスを機器の入出力インピーダンスに一致させるのではなく、低い出力インピーダンスで高い入力インピーダンスを駆動する。この場合には、ケーブルの静電容量、導体抵抗、シールド性能などが重要になる。

また、スピーカーの$4\,\Omega$や$8\,\Omega$はスピーカーの公称負荷インピーダンスであり、スピーカーケーブルの特性インピーダンスではない。スピーカーケーブルでは、特性インピーダンスをスピーカーに合わせることよりも、導体抵抗を十分小さくすることが求められる。

これに対して、S/PDIFやAES3などのデジタルオーディオ接続では、高速なパルスを伝送するため、$75\,\Omega$または$110\,\Omega$といった規定の特性インピーダンスが重要になる。

以上より、$\Omega$表示の意味は接続の種類によって異なり、アンテナ線やデジタル信号線では反射を防ぐための特性インピーダンス、アナログオーディオ機器では電圧伝送を決める入出力インピーダンス、スピーカーではアンプが駆動する負荷のインピーダンスを表していることが示された。

\end{document}
